% =============================================================================
%  RER DELAYS DATABASE PROJECT — REPORT
%  DBMS — April 2026
% =============================================================================

\documentclass[a4paper,11pt]{article}

% ── Packages ──────────────────────────────────────────────────────────
\usepackage[utf8]{inputenc}
\usepackage[T1]{fontenc}
\usepackage[english]{babel}
\usepackage{geometry}
\geometry{margin=2.5cm}
\usepackage{booktabs}
\usepackage{longtable}
\usepackage{listings}
\usepackage{xcolor}
\usepackage{hyperref}
\usepackage{caption}
\usepackage{float}
\usepackage{parskip}
\usepackage{amssymb}
\usepackage{amsmath}
\usepackage{graphicx}
\usepackage{tikz}
\usetikzlibrary{positioning, arrows.meta, calc, shapes.geometric, fit, backgrounds}
\usepackage[colorinlistoftodos, textsize=small]{todonotes}
\usepackage{mdframed}

% ── Placeholder and author-note commands ─────────────────────────────
\newcommand{\placeholder}[1]{%
  \begin{mdframed}[backgroundcolor=gray!15,
                   linecolor=gray!50,
                   linewidth=0.5pt,
                   innertopmargin=6pt,
                   innerbottommargin=6pt]
  \textit{\textcolor{gray}{#1}}
  \end{mdframed}%
}

\newcommand{\todoM}[1]{\todo[color=red!25,   inline]{\textbf{[M]} #1}}
\newcommand{\todoN}[1]{\todo[color=blue!20,  inline]{\textbf{[N]} #1}}
\newcommand{\todoL}[1]{\todo[color=green!25, inline]{\textbf{[L]} #1}}

% ── Code listing styles ───────────────────────────────────────────────
\definecolor{codebg}{gray}{0.96}
\definecolor{keyword}{RGB}{0,0,180}
\definecolor{comment}{RGB}{100,100,100}
\lstdefinestyle{sql}{
  language=SQL,
  basicstyle=\ttfamily\small,
  backgroundcolor=\color{codebg},
  keywordstyle=\color{keyword}\bfseries,
  commentstyle=\color{comment}\itshape,
  breaklines=true, frame=single, rulecolor=\color{gray},
  showstringspaces=false, tabsize=2, captionpos=b,
  morekeywords={TIMESTAMPTZ,DOUBLE,PRECISION,REFERENCES,TEXT,INTEGER,
                ILIKE,EXTRACT,ROUND,PERCENT_RANK,PARTITION,OVER,
                CASE,WHEN,THEN,ELSE,END,WIDTH_BUCKET,STDDEV,TRUNC,
                DATE_TRUNC,FILTER,HAVING,ISNULL}
}
\lstdefinestyle{python}{
  language=Python,
  basicstyle=\ttfamily\small,
  backgroundcolor=\color{codebg},
  keywordstyle=\color{keyword}\bfseries,
  commentstyle=\color{comment}\itshape,
  breaklines=true, frame=single, rulecolor=\color{gray},
  showstringspaces=false, tabsize=4, captionpos=b,
}
\lstdefinestyle{bash}{
  language=bash,
  basicstyle=\ttfamily\small,
  backgroundcolor=\color{codebg},
  frame=single, rulecolor=\color{gray},
  breaklines=true, showstringspaces=false,
}
% Style for EXPLAIN ANALYZE plan trees (no keyword highlighting)
\lstdefinestyle{explain}{
  basicstyle=\ttfamily\footnotesize,
  backgroundcolor=\color{codebg},
  frame=single, rulecolor=\color{gray},
  breaklines=true, captionpos=b,
}

% ── Title ─────────────────────────────────────────────────────────────
\title{
  \textbf{Analysis of RER Train Delays} \\[0.4em]
  \large A Relational Database Approach to IDFM Open Data \\[1.5em]
  \normalsize DBMS
}
\author{
  Nicola Cuzzola \and
  Matteo Faverio \and
  Lorenzo Marzuillo
}
\date{April 2026}

% =============================================================================
\begin{document}
% =============================================================================

\maketitle

% ── Abstract ─────────────────────────────────────────────────────────
\begin{abstract}
We design and implement a relational database for the analysis of train delays on the
Paris RER network, using a five-month corpus of real-time data self-collected from
the \^{I}le-de-France Mobilit\'{e}s (IDFM) PRIM API.
Starting from 150 raw daily CSV files containing 11,815,213 observations, we build
a normalised PostgreSQL schema following a star-snowflake design: one central fact
table with five dimension tables that correctly model the three-level stop hierarchy
of the IDFM network.
We describe the full ETL pipeline, the indexing strategy, and a catalogue of four
analytical queries alongside their \texttt{EXPLAIN ANALYZE} plans.
A controlled forced-plan experiment demonstrates the query planner's accuracy by
producing a $289\times$ slowdown when an index is forced on a full-table scan.
Finally, an interactive Streamlit dashboard exposes the four queries to non-technical
users.
\end{abstract}

\setcounter{tocdepth}{2}
\tableofcontents
\newpage


% =============================================================================
%  SECTION 1 — INTRODUCTION                                          [M]
% =============================================================================
\section{Introduction}
\label{sec:intro}

\subsection{Project context}

Paris's five RER lines form a cross-regional rail network carrying over three million
passengers daily.
Persistent delays are a well-documented quality-of-service problem, but no
historical archive of realised RER delays is available as a ready-to-use dataset.
\^{I}le-de-France Mobilit\'{e}s (IDFM) --- the regional transport authority ---
does provide the PRIM real-time API, which broadcasts current predicted passages
every few minutes for every stop on the network.
We therefore built our own delay corpus by polling that API continuously from a
dedicated virtual machine over five months, producing 150 daily CSV files and
11.8\,M aggregate observations.

This project uses the self-collected delay corpus to demonstrate the complete lifecycle of a
relational database system: from raw data ingestion through normalised schema design,
query optimisation, and interactive visualisation.
The system runs on a local PostgreSQL instance and is not a production deployment, but
the design principles and performance findings are representative of real database
engineering practice.

\subsection{Data source and collection pipeline}

The IDFM PRIM API exposes the endpoint
\texttt{/marketplace/estimated-timetable}, which returns a SIRI
\texttt{EstimatedTimetableDelivery} JSON snapshot of predicted passages at all stops
on demand.
Access requires an \texttt{apiKey} HTTP header.
We deployed a Python poller on a Linux virtual machine (Ubuntu\,24.04\,LTS),
orchestrated by a \texttt{systemd} timer.
The timer fires daily at 05:00 Europe/Paris and keeps the poller running until 02:30
the following morning.
At each polling instant (${\sim}300\,\text{s}$ cadence; observed 240–360\,s), the
poller downloads a full snapshot, aggregates the returned passages into one summary
row per \texttt{(stop\_id, line\_code)} pair, and appends the row to a daily CSV
file.
Because the collection window crosses midnight, polls between 00:00 and 02:29 are
stored under the \emph{previous} calendar date to preserve service-day boundaries.

\begin{table}[H]
\centering
\small
\begin{tabular}{ll}
\toprule
\textbf{Parameter} & \textbf{Value} \\
\midrule
Operating system    & Ubuntu 24.04 LTS \\
Scheduler           & \texttt{systemd} timer + service \\
Polling window      & 05:00–02:30 (Europe/Paris) \\
Target cadence      & 300\,s \\
Observed cadence    & 240–360\,s (occasional longer gaps on restart) \\
API endpoint        & \texttt{/marketplace/estimated-timetable} (SIRI format) \\
Authentication      & \texttt{apiKey} HTTP header \\
\bottomrule
\end{tabular}
\caption{VM data-collection setup.}
\label{tab:vm-setup}
\end{table}

After 150 service days of continuous polling (2025-11-12 to 2026-04-10),
the corpus contains \textbf{11,815,213 observations} across the five RER lines
(A, B, C, D, E).
Each row records aggregate statistics for one stop on one line during one polling
interval; the 9 columns are described in Table~\ref{tab:columns}.
Figure~\ref{fig:collection-pipeline} illustrates the complete pipeline from the
IDF network to the PostgreSQL database.

\begin{table}[H]
\centering
\small
\begin{tabular}{llp{7.0cm}}
\toprule
\textbf{Column} & \textbf{Type} & \textbf{Description} \\
\midrule
\texttt{poll\_at\_utc}    & TIMESTAMPTZ & Polling timestamp in UTC \\
\texttt{poll\_at\_local}  & TIMESTAMPTZ & Polling timestamp in Paris local time \\
\texttt{stop\_id}         & TEXT        & IDFM stop URI (three formats: Q\,/\,SP\,/\,BP) \\
\texttt{line\_code}       & TEXT        & RER line (e.g.\ \texttt{RER A}) \\
\texttt{mean\_delay\_s}   & DOUBLE      & Mean signed delay in seconds (negative = early) \\
\texttt{mean\_lateness\_s}& DOUBLE      & Mean lateness in seconds ($\geq 0$ always) \\
\texttt{n}                & INTEGER     & Passages contributing to this aggregate row \\
\texttt{n\_neg}           & INTEGER     & Passages with negative delay (early) \\
\texttt{n\_pos}           & INTEGER     & Passages with positive delay (late) \\
\bottomrule
\end{tabular}
\caption{CSV column schema (9 columns, identical across all 150 files).}
\label{tab:columns}
\end{table}

\begin{figure}[H]
  \centering
  \input{figures/collection_pipeline.tex}
  \caption{End-to-end pipeline: the PRIM API is polled every ${\sim}$5\,min from a
    Linux VM; each snapshot is aggregated per (stop, line) and appended to a daily
    CSV file; after 150 service days the CSVs are loaded into PostgreSQL via the
    ETL pipeline (\texttt{etl/load.py}).}
  \label{fig:collection-pipeline}
\end{figure}

\begin{figure}[H]
  \centering
  \includegraphics[width=0.82\textwidth]{figures/dataset_overview.pdf}
  \caption{Total polling observations per RER line across the 5-month corpus.
    RER\,A accounts for 30.7\% of observations, reflecting its higher stop density.}
  \label{fig:dataset-overview}
\end{figure}

\subsection{Goals and deliverables}

The project has four deliverables, each building on the previous:
\begin{enumerate}
  \item \textbf{Normalised PostgreSQL schema.}
    A star-snowflake schema with one fact table and five dimension tables,
    designed from first principles to avoid double-counting in the IDFM stop
    hierarchy.
  \item \textbf{Reproducible ETL pipeline.}
    A Python script (\texttt{etl/load.py}) that ingests all CSV and reference
    files, transforms them into the normalised schema, and is idempotent
    (safe to re-run without duplicating data).
  \item \textbf{Analytical query catalogue with \texttt{EXPLAIN ANALYZE}.}
    Four SQL queries covering the full spectrum of planner behaviours, each
    documented with its execution plan, selectivity, and index usage.
    The indexing strategy and a forced-plan degradation experiment are the
    core academic deliverable of this course.
  \item \textbf{Interactive Streamlit dashboard.}
    A web application exposing the four queries to non-technical users,
    demonstrating how index-aware query design translates into a
    responsive user experience.
\end{enumerate}

\subsection{Report structure}

Section~\ref{sec:profiling} describes the automated profiling of the raw CSV data
and its key findings.
Section~\ref{sec:conceptual} presents the conceptual model and entity-relationship
diagram.
Section~\ref{sec:schema} defines the physical PostgreSQL schema with full DDL.
Section~\ref{sec:etl} describes the ETL pipeline.
Section~\ref{sec:indexing} analyses the indexing strategy, and
Section~\ref{sec:queries} documents the four analytical queries with their
\texttt{EXPLAIN ANALYZE} plans and sample outputs.
Section~\ref{sec:optimisation} presents a controlled forced-plan degradation
experiment and the selectivity threshold model.
Section~\ref{sec:dashboard} describes the interactive dashboard.
Section~\ref{sec:conclusions} summarises findings, limitations, and future work.


\newpage
% =============================================================================
%  SECTION 2 — DATA PROFILING                                        [M]
% =============================================================================
\section{Data profiling}
\label{sec:profiling}

\subsection{Profiling methodology}

Before designing the schema we ran \texttt{etl/profile\_dataset.py}, an automated
profiler, across all 150 self-collected daily CSV files.
The script performs seven checks on each file individually and three additional
cross-file checks on the combined corpus, storing the full results in
\texttt{docs/profile\_summary.json} (regenerated on demand).

The six per-file checks are:
\begin{enumerate}
  \item \textbf{Schema consistency} --- all 9 column names and their order are
    identical across files.
  \item \textbf{Key uniqueness} --- no duplicate rows on the candidate key
    \texttt{(poll\_at\_utc, stop\_id, line\_code)} within a single file.
  \item \textbf{Null check} --- no nulls in any column.
  \item \textbf{$n$ constraint} --- \texttt{n\_neg + n\_pos $\leq$ n} holds for
    every row.
  \item \textbf{UTC$\to$local mapping} --- each UTC timestamp maps to exactly one
    local timestamp within the file.
  \item \textbf{Float type check} --- \texttt{mean\_delay\_s} and
    \texttt{mean\_lateness\_s} contain non-integer values (confirming
    \texttt{DOUBLE PRECISION} is the correct SQL type).
\end{enumerate}
The three cross-file checks verify key uniqueness globally, collect distinct line
codes, and confirm the overall date range.

\subsection{Results}

\begin{table}[H]
\centering
\begin{tabular}{llp{6cm}}
\toprule
\textbf{Check} & \textbf{Result} & \textbf{Notes} \\
\midrule
Schema consistency      & \checkmark Pass & Identical 9-column layout across all 150 files \\
Key uniqueness          & \checkmark Pass & 0 duplicate (poll\_at\_utc, stop\_id, line\_code) \\
Null check              & \checkmark Pass & No nulls in any column in any file \\
Constraint: $n_{\text{neg}} + n_{\text{pos}} \leq n$ & \checkmark Pass & Holds globally \\
UTC $\to$ local mapping & \checkmark Pass & Deterministic within each file \\
Float types             & \checkmark Pass & mean\_delay\_s and mean\_late\_s are DOUBLE PRECISION \\
\bottomrule
\end{tabular}
\label{tab:profiling}
\end{table}

\subsection{Key observations}

Two findings from profiling directly shaped the schema design.

\textbf{Natural key uniqueness.}
The global duplicate-key count is zero: the triple
\texttt{(poll\_at\_utc, stop\_id, line\_code)} uniquely identifies every row across
all 11.8\,M observations.
This result directly justifies using it as the composite
\texttt{PRIMARY KEY} of the fact table (Section~\ref{sec:schema}), without needing
a surrogate key.

\textbf{Three-format \texttt{stop\_id}.}
Inspection of the delay data revealed three distinct URI patterns in the
\texttt{stop\_id} column: \texttt{StopPoint:Q:*} (quay-level identifiers),
\texttt{StopArea:SP:*} (monomodal stop areas), and \texttt{StopPoint:BP:*}
(boarding points).
These are not interchangeable: they belong to different levels of the IDFM stop
hierarchy. This observation drove the snowflake extension to the schema
(Section~\ref{sec:conceptual}).


\newpage
% =============================================================================
%  SECTION 3 — CONCEPTUAL MODEL                                      [M]
% =============================================================================
\section{Conceptual model}
\label{sec:conceptual}

\subsection{Entity-relationship diagram}

The \texttt{core} schema follows a \emph{star-snowflake} design: a single central
fact table joined to five dimension tables.
Three of those dimensions form a snowflake chain that encodes the IDFM stop hierarchy
(Section~\ref{sec:stop-hierarchy}).
Figure~\ref{fig:er} shows the complete entity-relationship diagram; all relationships
are many-to-one (M:1) from the fact table or from a child dimension to its parent.

\begin{figure}[H]
  \centering
  \input{figures/er_diagram.tex}
  \caption{Entity-relationship diagram of the \texttt{core} schema.
    PK columns are underlined; FK columns note their target table.
    The snowflake chain \texttt{dim\_stop} $\to$ \texttt{dim\_monomodal\_stop}
    $\to$ \texttt{dim\_station} models the three-level IDFM stop hierarchy.}
  \label{fig:er}
\end{figure}

\subsection{The stop hierarchy}
\label{sec:stop-hierarchy}

A critical finding from data profiling is that \texttt{stop\_id} is \emph{not}
a station identifier: it refers to a specific level in the IDFM network hierarchy.
Three URI patterns appear in the delay data:

\begin{itemize}
  \item \textbf{Q} (\texttt{StopPoint:Q:*}) --- a quay or individual platform face.
    Each physical station may have several quay entries.
  \item \textbf{SP} (\texttt{StopArea:SP:*}) --- a monomodal stop area, grouping
    all quays of the same transport mode at a station.
  \item \textbf{BP} (\texttt{StopPoint:BP:*}) --- a boarding point (less common;
    39 entries in the dataset).
\end{itemize}

The reference file \texttt{data/stations.csv} exposes the three-level hierarchy:

\[
\scalebox{0.88}{$\displaystyle
  \underbrace{\texttt{dim\_station}}_{\text{475 stations}}
  \;\leftarrow\;
  \underbrace{\texttt{dim\_monomodal\_stop}}_{\text{476 modal groups}}
  \;\leftarrow\;
  \underbrace{\texttt{dim\_stop}}_{\text{908 stop entries}}
  \;\leftarrow\;
  \underbrace{\texttt{fact\_delay\_events}}_{\text{11.8\,M rows}}
$}
\]

\textbf{Why SP rows are used for station-level aggregation.}
Each physical station is represented in the delay data by \emph{both} Q-type rows
(one per platform) and an SP-type row (representing the whole modal group).
If both were included in a station-level aggregation (e.g.\ Q7, the geographic
map), each station's trains would be counted multiple times.
Filtering \texttt{WHERE s.stop\_type = 'SP'} in analytical queries ensures each
observation is counted exactly once at the station level.

Note: 390 of the 570 Q-type stop identifiers from the first file could not be
resolved in \texttt{stations.csv} (their \texttt{monomodal\_code} is \texttt{NULL}).
This is expected --- the reference file covers a curated subset of the full network
--- and is not treated as a data error.

\subsection{Normalisation}

The schema is in Boyce-Codd Normal Form (BCNF).
A relation is in BCNF if and only if, for every non-trivial functional dependency
$X \to Y$, $X$ is a superkey.
We verify this for each table.

\begin{itemize}
  \item \textbf{\texttt{fact\_delay\_events}.}
    Candidate key: \texttt{(poll\_at\_utc, stop\_id, line\_code)}.
    The five metric columns (\texttt{mean\_delay\_s}, \texttt{mean\_lateness\_s},
    \texttt{n}, \texttt{n\_neg}, \texttt{n\_pos}) are direct measurements of that
    specific polling event; no subset of the key determines them independently.
    BCNF. \\
  \item \textbf{\texttt{dim\_time}.}
    Candidate key: \texttt{poll\_at\_utc}.
    FD: \texttt{poll\_at\_utc} $\to$ \texttt{poll\_at\_local}.
    The determinant is the PK. BCNF. \\
  \item \textbf{\texttt{dim\_line}.}
    Candidate key: \texttt{line\_code}. No other attributes. Trivially BCNF. \\
  \item \textbf{\texttt{dim\_stop}.}
    Candidate key: \texttt{stop\_id}.\\
    FD: \texttt{stop\_id} $\to$ \texttt{\{stop\_type, quay\_code, monomodal\_code\}}.
    Determinant is the PK. BCNF. \\
  \item \textbf{\texttt{dim\_monomodal\_stop}.}
    Candidate key: \texttt{monomodal\_code}.\\
    FD: \texttt{monomodal\_code} $\to$ \texttt{\{monomodal\_stop\_id, station\_id\}}.
    BCNF. \\
  \item \textbf{\texttt{dim\_station}.}
    Candidate key: \texttt{station\_id}.
    FD: \texttt{station\_id} $\to$ \texttt{\{stop\_name, stop\_lat, stop\_lon,
    zone\_id\}}.
    BCNF.
\end{itemize}

\textbf{Deliberate design choices.}
\texttt{line\_code} appears as a column in the fact table rather than being
derived through the dimension chain.
This is not a denormalisation: it is part of the natural composite key of the
observation (a polling event is inherently line-specific) and its presence avoids
an extra join for every line-filtered query.
Similarly, \texttt{poll\_at\_local} is materialised in \texttt{dim\_time} rather
than computed at query time, since the UTC$\to$local conversion must account for
DST transitions and is non-trivial to reproduce correctly in SQL.


\newpage
% =============================================================================
%  SECTION 4 — PHYSICAL SCHEMA                                       [M]
% =============================================================================
\section{Physical schema}
\label{sec:schema}

\subsection{Schema overview}

\begin{table}[H]
\centering
\begin{tabular}{llr}
\toprule
\textbf{Table} & \textbf{Role} & \textbf{Rows} \\
\midrule
\texttt{fact\_delay\_events} & Central fact table (the measurements) & 11,815,213 \\
\texttt{dim\_time}           & Polling timestamps (UTC + local)      & 37,146 \\
\texttt{dim\_line}           & RER line codes                        & 5 \\
\texttt{dim\_station}        & Physical stations (name, coordinates) & 475 \\
\texttt{dim\_monomodal\_stop}& Monomodal stop areas                  & 476 \\
\texttt{dim\_stop}           & Stop-level entries (Q / SP / BP)      & 908 \\
\bottomrule
\end{tabular}
\caption{Tables in the \texttt{core} schema.}
\label{tab:tables}
\end{table}

\subsection{Fact table}

The fact table is the centrepiece of the schema.
Its \texttt{PRIMARY KEY} uses the natural composite key identified during profiling
--- zero global duplicates confirmed this choice.
Two \texttt{CHECK} constraints enforce the domain invariants validated by the
profiler: the train-count decomposition ($n_{\text{neg}} + n_{\text{pos}} \leq n$)
and the non-negativity of \texttt{mean\_lateness\_s}.
Both metric columns use \texttt{DOUBLE PRECISION}, confirmed by the profiler's
float-type check.
All three \texttt{REFERENCES} clauses enforce referential integrity to the
dimension tables.

\begin{lstlisting}[style=sql, caption={Fact table DDL (\texttt{sql/schema/02\_core.sql}).}]
CREATE TABLE core.fact_delay_events (
    poll_at_utc      TIMESTAMPTZ      NOT NULL
                     REFERENCES core.dim_time(poll_at_utc),
    stop_id          TEXT             NOT NULL
                     REFERENCES core.dim_stop(stop_id),
    line_code        TEXT             NOT NULL
                     REFERENCES core.dim_line(line_code),
    mean_delay_s     DOUBLE PRECISION NOT NULL,
    mean_lateness_s  DOUBLE PRECISION NOT NULL,
    n                INTEGER          NOT NULL CHECK (n >= 0),
    n_neg            INTEGER          NOT NULL CHECK (n_neg >= 0),
    n_pos            INTEGER          NOT NULL CHECK (n_pos >= 0),
    PRIMARY KEY (poll_at_utc, stop_id, line_code),
    CHECK (n_neg + n_pos <= n),
    CHECK (mean_lateness_s >= 0)
);
\end{lstlisting}

\subsection{Dimension tables}

The five dimension tables are straightforward; Table~\ref{tab:tables} lists their
row counts.
Notable design decisions:
\begin{itemize}
  \item \texttt{dim\_station.zone\_id} and \texttt{station\_code} are nullable
    because they are absent for a subset of quays in the source file.
  \item \texttt{dim\_stop.stop\_type} has a \texttt{CHECK} constraint restricting
    values to \texttt{'Q'}, \texttt{'SP'}, or \texttt{'BP'}.
  \item \texttt{dim\_stop.monomodal\_code} is nullable: it is \texttt{NULL} for
    Q-type stops not resolvable in the source file \texttt{stations.csv} (390 entries) and for
    all BP-type stops.
    This is expected and documented, not a data error.
\end{itemize}

\begin{lstlisting}[style=sql, caption={Dimension table DDL (\texttt{sql/schema/02\_core.sql}).}]
CREATE TABLE core.dim_time (
    poll_at_utc   TIMESTAMPTZ PRIMARY KEY,
    poll_at_local TIMESTAMPTZ NOT NULL
);

CREATE TABLE core.dim_line (
    line_code TEXT PRIMARY KEY
);

CREATE TABLE core.dim_station (
    station_id    INTEGER          PRIMARY KEY,
    stop_name     TEXT             NOT NULL,
    stop_lat      DOUBLE PRECISION NOT NULL,
    stop_lon      DOUBLE PRECISION NOT NULL,
    zone_id       TEXT,          -- nullable: absent for ~10% of quays
    station_code  TEXT           -- nullable: absent for ~43% of quays
);

CREATE TABLE core.dim_monomodal_stop (
    monomodal_code    INTEGER PRIMARY KEY,
    monomodal_stop_id TEXT    NOT NULL,
    station_id        INTEGER NOT NULL
                      REFERENCES core.dim_station(station_id)
);

CREATE TABLE core.dim_stop (
    stop_id         TEXT    PRIMARY KEY,
    stop_type       TEXT    CHECK (stop_type IN ('Q', 'SP', 'BP')),
    quay_code       INTEGER,
    monomodal_code  INTEGER
                    REFERENCES core.dim_monomodal_stop(monomodal_code)
    -- NULL for unresolved Q-type and all BP-type stops
);
\end{lstlisting}

\subsection{Staging tables}

The \texttt{staging} schema provides an immutable landing zone for all raw data.
No transformation is applied to staging tables; the raw data is preserved exactly
as it arrives.

\texttt{staging.delay\_events} mirrors the CSV column layout exactly, with two
additional audit columns: \texttt{source\_file} (the originating filename, used to
detect already-loaded files) and \texttt{loaded\_at} (the ingestion timestamp).

\texttt{staging.stops} is the raw load of \texttt{data/stations.csv}, one row per
quay (2,533 rows), with all 11 source columns preserved verbatim.
It encodes the three-level stop hierarchy in denormalised form; the transformation
step (Section~\ref{sec:etl}) normalises it into \texttt{dim\_station},
\texttt{dim\_monomodal\_stop}, and the enrichment of \texttt{dim\_stop}.

The DDL for both staging tables is in \texttt{sql/schema/01\_staging.sql}.


\newpage
% =============================================================================
%  SECTION 5 — ETL PIPELINE                                          [M]
% =============================================================================
\section{ETL pipeline}
\label{sec:etl}

\subsection{Architecture}

The pipeline (\texttt{etl/load.py}) follows a two-schema architecture with two
parallel data streams, illustrated in Figure~\ref{fig:etl-flow}.

\begin{itemize}
  \item \textbf{Delay stream:}\\
    \texttt{data/raw/*.csv} (150 files) $\to$ \texttt{staging.delay\_events}
    $\to$ \texttt{core.fact\_delay\_events} + \texttt{dim\_time}, \texttt{dim\_line},
    \texttt{dim\_stop}.
  \item \textbf{Station stream:}\\
    \texttt{data/stations.csv} $\to$ \texttt{staging.stops}
    $\to$ \texttt{core.dim\_station} + \texttt{dim\_monomodal\_stop};
    also enriches \texttt{dim\_stop} with \texttt{stop\_type} and
    \texttt{monomodal\_code}.
\end{itemize}

The staging layer exists to keep raw data immutable.
If the core transform needs to be rerun (e.g.\ after a schema change), the
CSVs do not need to be reloaded: the staging tables remain intact.

\begin{figure}[H]
  \centering
  \input{figures/etl_flow.tex}
  \caption{ETL pipeline data flow.
    Orange arrows: bulk \texttt{COPY FROM STDIN}.
    Blue arrows: \texttt{INSERT \ldots{} SELECT \ldots{} ON CONFLICT DO NOTHING}.
    Dashed green arrow: \texttt{dim\_stop} enrichment from \texttt{staging.stops}.}
  \label{fig:etl-flow}
\end{figure}

\subsection{Loading steps}

The pipeline runs in five phases (as reported by \texttt{load.py}):

\begin{enumerate}
  \item \textbf{Schema setup.}
    Creates the \texttt{staging} and \texttt{core} schemas and all tables if they
    do not already exist (\texttt{01\_staging.sql}, \texttt{02\_core.sql}).
    The optional \texttt{--reset} flag drops all schemas first, enabling a clean
    rebuild for development.

  \item \textbf{Bulk load of delay CSVs.}
    Each of the 150 files is loaded into \texttt{staging.delay\_events} using
    \texttt{COPY FROM STDIN} (psycopg's binary protocol).
    \texttt{COPY} bypasses row-by-row overhead, making it 10--100$\times$ faster
    than \texttt{INSERT} for large files.
    Before copying, the script checks whether the file is already recorded in
    \texttt{source\_file}; if so, it is skipped (idempotency mechanism).

  \item \textbf{Load station reference file.}
    \texttt{data/stations.csv} is loaded into \texttt{staging.stops} via
    \texttt{COPY FROM STDIN}.
    Skipped if \texttt{staging.stops} already contains rows.

  \item \textbf{Transform delay data $\to$ core.}
    Executes \texttt{03\_load\_core\_from\_staging.sql}, which populates
    \texttt{dim\_time}, \texttt{dim\_line}, \texttt{dim\_stop}, and
    \texttt{fact\_delay\_events} via \texttt{INSERT \ldots{} SELECT \ldots{} ON
    CONFLICT DO NOTHING}.
    Indexes are created \emph{after} this bulk insert, not before, to avoid the
    per-row B-tree update overhead during the initial load.

  \item \textbf{Transform stop data $\to$ core.}
    Executes \texttt{04\_load\_stops.sql}, which populates \texttt{dim\_station}
    and \texttt{dim\_monomodal\_stop}, and then also enriches \texttt{dim\_stop} with
    \texttt{stop\_type}, \texttt{quay\_code}, and \texttt{monomodal\_code}.
    The 390 Q-type stops absent from \texttt{stations.csv} retain
    \texttt{monomodal\_code = NULL}; (documented, not RER stops).
\end{enumerate}

\subsection{Idempotency}

The pipeline is \emph{idempotent}: re-running \texttt{etl/load.py} on an already-loaded
database produces no duplicate data and no errors.
This property is achieved through two complementary mechanisms:

\begin{itemize}
  \item \textbf{\texttt{ON CONFLICT DO NOTHING}} on every
    \texttt{INSERT \ldots{} SELECT} in the transform SQL files.
    Any row whose primary key already exists in the target table is silently
    skipped, leaving existing data unchanged.
  \item \textbf{Source-file tracking.}
    Before copying a CSV, \texttt{load.py} queries
    \texttt{staging.delay\_events} for rows with that \texttt{source\_file}.
    If any exist, the file is entirely skipped, avoiding redundant staging data.
\end{itemize}

For a full rebuild (e.g.\ after a schema migration), passing \texttt{--reset}
drops all schemas before running, effectively converting the pipeline into a
destructive bulk reload.
This flag is intended for development only.


\newpage
% =============================================================================
%  SECTION 6 — INDEXING STRATEGY                                     [N]
% =============================================================================
\section{Indexing strategy}
\label{sec:indexing}

\subsection{Why indexes matter}

A B-tree index operates much like the alphabetical index at the back of a book. Instead of reading every page to find a specific topic, the database looks up the ``page number'' in the index and jumps directly to it. Our main fact table contains 11.8 million rows, occupying approximately 1\,GB on disk. Without indexes, every query would be forced to perform a full sequential scan, reading the entire table from start to finish. However, there is a fundamental trade-off: while indexes drastically speed up selective reads, they slow down bulk write operations because the index tree must be updated for every new row. For this reason, we deliberately created the indexes \emph{after} the bulk data load.

\subsection{Indexes created}

\begin{table}[H]
\centering
\small
\resizebox{\linewidth}{!}{%
\begin{tabular}{lllp{5cm}}
\toprule
\textbf{Name} & \textbf{Columns} & \textbf{Size} & \textbf{Purpose} \\
\midrule
\texttt{idx\_fact\_line\_time} & \texttt{(line\_code, poll\_at\_utc)} & 82 MB &
  Queries filtering by line + date range (Q6, Q7) \\
\texttt{idx\_fact\_time} & \texttt{(poll\_at\_utc)} & 81 MB &
  Queries filtering by date only \\
\texttt{idx\_fact\_stop} & \texttt{(stop\_id)} & 81 MB &
  Single-station lookups (Q5) \\
\texttt{idx\_stop\_monomodal} & \texttt{(monomodal\_code)} partial & 16\,kB &
  Joining \texttt{dim\_stop} to \texttt{dim\_monomodal\_stop} \\
\midrule
Primary key & \texttt{(poll\_at\_utc, stop\_id, line\_code)} & 863 MB &
  Uniqueness enforcement \\
\bottomrule
\end{tabular}%
}% end resizebox
\caption{B-tree indexes on the \texttt{core} schema.}
\label{tab:indexes}
\end{table}

We designed four specific indexes tailored to our analytical queries:
\begin{itemize}
  \item \textbf{\texttt{idx\_fact\_line\_time}}: Designed for queries filtering by both line and date range (Q6, Q7). The column order is critical: \texttt{line\_code} is the leading column because it handles equality predicates (e.g.\ \texttt{line = 'RER A'}), allowing the engine to jump directly to that section of the B-tree. The second column, \texttt{poll\_at\_utc}, is then used to scan the date range within that specific line. This follows the standard composite index design rule: \emph{equality predicate first, range predicate second}.
  \item \textbf{\texttt{idx\_fact\_time}}: Designed for queries filtering by date across all lines, without a line equality predicate.
  \item \textbf{\texttt{idx\_fact\_stop}}: Designed for single-station point lookups (Q5), allowing near-instant retrieval of all rows associated with a specific \texttt{stop\_id}.
  \item \textbf{\texttt{idx\_stop\_monomodal}}: A \emph{partial index}, meaning it only indexes rows matching the condition \texttt{WHERE monomodal\_code IS NOT NULL}. Because it only stores rows that actually belong to a monomodal grouping, it is extremely compact (16\,kB) and accelerates hierarchical joins between \texttt{dim\_stop} and \texttt{dim\_monomodal\_stop}.
\end{itemize}

\subsection{The query planner and selectivity}

The PostgreSQL query planner dynamically chooses the execution plan based on \emph{selectivity}---the fraction of the table a query needs to retrieve. The planner calculates the expected cost and follows a strict decision rule:
\begin{itemize}
  \item \textbf{$<{\sim}10$--$15\%$ selectivity $\Rightarrow$ Index Scan:} When requesting a small fraction of data, it is cost-effective to perform \emph{random I/O} (accessing specific disk pages out of order via the index) because only a few pages need to be fetched.
  \item \textbf{$>{\sim}15$--$20\%$ selectivity $\Rightarrow$ Sequential Scan:} When requesting a large portion of the table, the planner ignores the index. Random I/O becomes extremely expensive because jumping between scattered pages incurs high disk latency. Instead, \emph{sequential I/O} is used: pages are read contiguously, allowing the OS read-ahead prefetcher to load upcoming pages into RAM before the database even requests them.
\end{itemize}
This fundamental concept of cost-based plan selection is explored further and practically demonstrated in Section~\ref{sec:optimisation}.


\newpage
% =============================================================================
%  SECTION 7 — QUERY CATALOGUE                                       [N]
% =============================================================================
\section{Query catalogue}
\label{sec:queries}

\noindent
We present 4 queries that together cover the full range of optimisation
behaviours: index point lookup (Q5), adaptive index/scan selection (Q6),
compound-predicate index (Q7), and full-table aggregation (Q8).
Each query is linked to a specific dashboard widget.


% ──────────────────────────────────────────────────────────────────────
\subsection{Q5 --- Station delay timeline}
% ──────────────────────────────────────────────────────────────────────

\textbf{Use case:} A commuter wants to see how delays at their station have evolved over 5 months.
\subsubsection{SQL}

\begin{lstlisting}[style=sql, caption={Q5 --- station delay timeline
  (\texttt{sql/queries/q5\_station\_delay\_timeline.sql}).}]
SELECT
    date_trunc('day', f.poll_at_utc)::date          AS day,
    round(avg(f.mean_delay_s)::numeric,    2)       AS avg_delay_s,
    round(avg(f.mean_lateness_s)::numeric, 2)       AS avg_lateness_s,
    round(stddev(f.mean_delay_s)::numeric, 2)       AS stddev_delay_s,
    sum(f.n)                                        AS total_trains,
    count(*)                                        AS n_polls
FROM core.fact_delay_events    f
JOIN core.dim_stop              s  ON s.stop_id        = f.stop_id
JOIN core.dim_monomodal_stop   ms ON ms.monomodal_code = s.monomodal_code
JOIN core.dim_station          ds ON ds.station_id     = ms.station_id
WHERE s.stop_type  = 'SP'
  AND ds.stop_name = :station_name     -- e.g. 'Chatelet-Les-Halles'
GROUP BY day
ORDER BY day;
\end{lstlisting}

\subsubsection{Relational algebra}

The query selects SP-type stops matching the given station name, joins through
the stop hierarchy, and aggregates to daily grain.
Let $p$ denote the station name parameter.
\begin{align*}
  &\gamma_{\text{day};\;
    \text{avg}(\texttt{mean\_delay\_s}),\,
    \text{avg}(\texttt{mean\_lateness\_s}),\,
    \text{sum}(\texttt{n}),\,
    \text{count}(*)} \\
  &\quad\Bigl(
    \sigma_{\texttt{stop\_name}=p\;\land\;\texttt{stop\_type}=\text{`SP'}}
    \bigl(
      \texttt{fact} \bowtie \texttt{dim\_stop} \bowtie
      \texttt{dim\_monomodal\_stop} \bowtie \texttt{dim\_station}
    \bigr)
  \Bigr)
\end{align*}

\subsubsection{EXPLAIN ANALYZE}

The execution takes 9,539\,ms on a cold cache and $<$1\,s on a warm cache.
The planner executes the following steps:
\begin{enumerate}
  \item \textbf{Seq Scan} on \texttt{dim\_station} (475 rows) to resolve the
    station name to its \texttt{station\_id}.
  \item \textbf{Hash Join} to \texttt{dim\_monomodal\_stop} (476 rows, fits in memory).
  \item \textbf{Index Scan} on \texttt{dim\_stop} via \texttt{idx\_stop\_monomodal}
    (partial index, 16\,kB).
  \item \textbf{Bitmap Index Scan} on the fact table using \texttt{idx\_fact\_stop}.
    This is the critical step: it isolates exactly 83,633 rows, representing
    \textbf{0.7\% selectivity} of the 11.8\,M-row table.
  \item \textbf{GroupAggregate} to compute the final daily averages.
\end{enumerate}
The index is chosen because 0.7\% selectivity is well below the ${\sim}15\%$
break-even threshold: fetching only 83,633 scattered pages via random I/O is
far cheaper than reading the entire 1\,GB table sequentially.

\subsubsection{Sample output}

\begin{table}[H]
\centering
\small
\resizebox{\linewidth}{!}{%
\begin{tabular}{lrrrrrr}
\toprule
\textbf{day} & \textbf{avg\_delay\_s} & \textbf{avg\_lateness\_s}
             & \textbf{stddev\_delay\_s} & \textbf{total\_trains} & \textbf{n\_polls} \\
\midrule
2025-11-12 &  45.15 &  47.52 &  63.95 &  461 & 154 \\
2025-11-13 & 157.44 & 160.13 & 267.04 &  621 & 200 \\
2025-11-14 &  66.09 &  68.02 &  88.07 &  601 & 207 \\
2025-11-15 &  58.32 &  60.08 &  90.26 &  514 & 203 \\
2025-11-16 &  50.87 &  52.44 & 111.48 &  390 & 167 \\
\bottomrule
\end{tabular}%
}% end resizebox
\caption{Q5 sample output for Ch\^{a}telet les Halles (first 5 days of corpus).}
\end{table}


% ──────────────────────────────────────────────────────────────────────
\subsection{Q6 --- Line comparison over a date range}
% ──────────────────────────────────────────────────────────────────────

\textbf{Use case:} An operator wants to compare all 5 lines for a chosen period.

\subsubsection{SQL}

\begin{lstlisting}[style=sql, caption={Q6 --- line comparison
  (\texttt{sql/queries/q6\_line\_comparison\_date\_range.sql}).}]
SELECT
    line_code,
    round(avg(mean_delay_s)::numeric,    2)        AS avg_delay_s,
    round(avg(mean_lateness_s)::numeric, 2)        AS avg_lateness_s,
    round(stddev(mean_delay_s)::numeric, 2)        AS stddev_delay_s,
    round(max(mean_delay_s)::numeric,    2)        AS max_delay_s,
    sum(n)                                         AS total_trains,
    count(*)                                       AS n_observations
FROM core.fact_delay_events
WHERE poll_at_utc >= :date_from          -- e.g. '2026-03-01'
  AND poll_at_utc <  :date_to            -- e.g. '2026-03-08'
GROUP BY line_code
ORDER BY avg_lateness_s DESC;
\end{lstlisting}

\subsubsection{Relational algebra}

No join is needed; the query aggregates the fact table directly.
Let $d_1$ and $d_2$ denote the start and end timestamps.
\begin{align*}
  &\gamma_{\texttt{line\_code};\;
    \text{avg}(\texttt{mean\_delay\_s}),\,\text{avg}(\texttt{mean\_lateness\_s}),\,
    \text{max}(\texttt{mean\_delay\_s}),\,\text{sum}(\texttt{n}),\,\text{count}(*)} \\
  &\quad\left(
    \sigma_{d_1 \leq \texttt{poll\_at\_utc} < d_2}
    \bigl(\texttt{fact\_delay\_events}\bigr)
  \right)
\end{align*}

\subsubsection{EXPLAIN ANALYZE --- plan changes with selectivity}

This query clearly demonstrates how the planner adapts to selectivity:
\begin{itemize}
  \item \textbf{Case 1 --- narrow window (7 days):} The query requests
    ${\sim}551{,}000$ rows (4.7\% of the table). The planner uses a
    \textbf{Bitmap Index Scan} on \texttt{idx\_fact\_line\_time},
    completing in 4,323\,ms.
  \item \textbf{Case 2 --- full corpus (5 months):} The query requests all
    11.8\,M rows (100\%). The planner completely ignores the index and switches
    to a \textbf{Parallel Seq Scan}, finishing in ${\sim}1{,}400$\,ms.
\end{itemize}
The crossover point lies around 10--15\% selectivity; beyond this threshold,
sequential reading surpasses index traversal in efficiency.

\begin{table}[H]
\centering
\begin{tabular}{lrrl}
\toprule
\textbf{Date window} & \textbf{Rows (\%)} & \textbf{Time (ms)} &
  \textbf{Plan chosen} \\
\midrule
7 days      & 551\,K (4.7\%)  & 4,323        & Bitmap Index Scan \\
Full corpus & 11.8\,M (100\%) & ${\sim}$1,400 & Parallel Seq Scan \\
\bottomrule
\end{tabular}
\caption{Q6: the planner chooses different plans for different selectivities.}
\label{tab:q6-plans}
\end{table}

\subsubsection{Sample output}

\begin{table}[H]
\centering
\small
\resizebox{\linewidth}{!}{%
\begin{tabular}{lrrrrrr}
\toprule
\textbf{line\_code} & \textbf{avg\_delay\_s} & \textbf{avg\_lateness\_s}
  & \textbf{stddev\_delay\_s} & \textbf{max\_delay\_s} & \textbf{total\_trains} & \textbf{n\_obs} \\
\midrule
RER C &  86.52 &  98.40 & 281.96 & 3596.00 & 190{,}734 & 120{,}968 \\
RER D &  78.20 &  86.61 & 230.47 & 3444.00 & 155{,}911 & 101{,}886 \\
RER B &  59.59 &  74.39 & 175.37 & 3585.00 & 212{,}330 & 108{,}999 \\
RER A &  58.42 &  74.22 & 186.97 & 3499.00 & 297{,}057 & 168{,}734 \\
RER E &  54.43 &  67.33 & 199.12 & 3270.00 &  88{,}454 &  50{,}845 \\
\bottomrule
\end{tabular}%
}% end resizebox
\caption{Q6 sample output for the 7-day window 2026-03-01 to 2026-03-08, ordered by average lateness.}
\end{table}


% ──────────────────────────────────────────────────────────────────────
\subsection{Q7 --- Geographic delay map}
% ──────────────────────────────────────────────────────────────────────

\textbf{Use case:} A data journalist wants to visualise which parts of
the RER A network suffer the most delays in March.

\subsubsection{SQL}

\begin{lstlisting}[style=sql, caption={Q7 --- geographic delay map
  (\texttt{sql/queries/q7\_geographic\_delay\_map.sql}).}]
SELECT
    ds.stop_name,
    ds.stop_lat,
    ds.stop_lon,
    ds.zone_id,
    f.line_code,
    round(avg(f.mean_delay_s)::numeric,    2)      AS avg_delay_s,
    round(avg(f.mean_lateness_s)::numeric, 2)      AS avg_lateness_s,
    sum(f.n)                                       AS total_trains,
    count(*)                                       AS n_polls
FROM core.fact_delay_events    f
JOIN core.dim_stop              s  ON s.stop_id        = f.stop_id
JOIN core.dim_monomodal_stop   ms ON ms.monomodal_code = s.monomodal_code
JOIN core.dim_station          ds ON ds.station_id     = ms.station_id
WHERE s.stop_type  = 'SP'
  AND f.line_code  = :line_code        -- e.g. 'RER A'
  AND f.poll_at_utc >= :date_from      -- e.g. '2026-03-01'
  AND f.poll_at_utc <  :date_to        -- e.g. '2026-04-01'
GROUP BY ds.stop_name, ds.stop_lat, ds.stop_lon, ds.zone_id, f.line_code
HAVING count(*) >= 10
ORDER BY avg_lateness_s DESC;
\end{lstlisting}

\subsubsection{Relational algebra}

The query combines the compound date+line predicate of Q6 with the
three-table join of Q5, and adds a \textsc{Having} filter after
aggregation. Let $\ell$ denote the line parameter and $[d_1, d_2)$ the date range.
\begin{align*}
  &\sigma_{\text{count}(*) \geq 10}
  \Bigl(
    \gamma_{\text{name, lat, lon, line};\;
      \text{avg}(\texttt{mean\_delay\_s}),\,\text{avg}(\texttt{mean\_lateness\_s}),\,\text{count}(*)} \\
  &\quad\bigl(
    \sigma_{\texttt{line\_code}=\ell \;\land\;
            d_1 \leq \texttt{poll\_at\_utc} < d_2 \;\land\;
            \texttt{stop\_type}=\text{`SP'}} \\
  &\qquad\left(
    \texttt{fact} \bowtie \texttt{dim\_stop} \bowtie
    \texttt{dim\_monomodal\_stop} \bowtie \texttt{dim\_station}
  \right)
  \bigr)
  \Bigr)
\end{align*}

\subsubsection{EXPLAIN ANALYZE}

For the RER\,A test case (March 2026), the query executed in 4,516\,ms
and returned 46 stations. The planner follows these steps:
\begin{enumerate}
  \item \textbf{Bitmap Index Scan} on \texttt{idx\_fact\_line\_time}.
    The compound predicate (\texttt{line\_code = 'RER A'} AND date range) is fully
    satisfied by the index because \texttt{line\_code} is the leading column
    (equality predicate) and \texttt{poll\_at\_utc} is the second column (range predicate).
  \item \textbf{Parallel Bitmap Heap Scan} on \texttt{fact\_delay\_events} to
    fetch the actual rows identified by step~1.
  \item \textbf{Hash Join $\times$ 3} against \texttt{dim\_stop} (908 rows),
    \texttt{dim\_monomodal\_stop} (476 rows), and \texttt{dim\_station} (475 rows).
    All three dimension tables fit entirely in RAM, making these joins negligible in cost.
  \item \textbf{GroupAggregate + HAVING} filter to produce the final 46 station rows.
\end{enumerate}

Had the index column order been reversed (\texttt{poll\_at\_utc} first, then
\texttt{line\_code}), the index would be far less efficient: the engine would
need to scan multiple disjoint time segments across all lines and then filter
by \texttt{line\_code} afterwards, erasing most of the index benefit.

\begin{lstlisting}[style=explain,
  caption={EXPLAIN (ANALYZE, BUFFERS) output for Q7: RER A, March 2026.
    Execution time: 4,020\,ms (cold cache).
    Key: \texttt{idx\_fact\_line\_time} is hit in 70\,ms; the three
    dimension-table hashes are each $<$1\,ms.},
  label=lst:q7-plan]
Sort  (actual rows=46, Execution Time=4020 ms)
  Sort Key: avg_lateness_s DESC
  ->  Finalize GroupAggregate  (actual rows=46)
        Filter: (count(*) >= 10)
        ->  Gather Merge  (Workers Launched: 2)
              ->  Partial HashAggregate  (actual rows=46 per worker)
                    ->  Hash Join  (dim_station hash, 475 rows, 44 kB)
                          ->  Hash Join  (dim_monomodal_stop hash, 476 rows, 27 kB)
                                ->  Hash Join  (dim_stop hash, 243 SP rows, 23 kB)
                                      ->  Parallel Bitmap Heap Scan
                                            on fact_delay_events
                                            (actual rows=244,160 per worker)
                                            Heap Blocks: exact=19315 lossy=10971
                                            ->  Bitmap Index Scan
                                                  on idx_fact_line_time
                                                  Index Cond: line_code='RER A'
                                                    AND poll_at_utc in [Mar 2026]
                                                  (rows=732,480  time=70 ms)
                                                  Buffers: shared read=645
                                      ->  Hash on dim_stop
                                            Filter: stop_type = 'SP'
                                ->  Hash on dim_monomodal_stop
                          ->  Hash on dim_station
Planning Time: 31 ms  |  Execution Time: 4,020 ms
\end{lstlisting}

\subsubsection{Sample output}

\begin{table}[H]
\centering
\small
\resizebox{\linewidth}{!}{%
\begin{tabular}{lrrrrrr}
\toprule
\textbf{stop\_name} & \textbf{avg\_delay\_s} & \textbf{avg\_lateness\_s}
  & \textbf{stop\_lat} & \textbf{stop\_lon} & \textbf{zone} & \textbf{n\_polls} \\
\midrule
Cergy Pr\'{e}fecture      & 76.91 & 78.79 & 49.0363 & 2.0798 & 5 & 5{,}590 \\
Conflans Fin d'Oise       & 75.67 & 76.52 & 48.9912 & 2.0746 & 5 & 5{,}971 \\
Neuville Universit\'{e}   & 75.59 & 76.02 & 49.0137 & 2.0783 & 5 & 5{,}737 \\
Ach\`{e}res Ville         & 73.72 & 75.11 & 48.9706 & 2.0777 & 5 & 6{,}049 \\
Ach\`{e}res Grand Cormier & 66.11 & 66.71 & 48.9554 & 2.0937 & 5 & 4{,}927 \\
\bottomrule
\end{tabular}%
}% end resizebox
\caption{Q7 sample output: top 5 stations by average lateness, RER A, March 2026.
  All are western-branch terminus stations, reflecting their distance from Paris.}
\end{table}


% ──────────────────────────────────────────────────────────────────────
\subsection{Q8 --- Peak vs.\ off-peak comparison}
% ──────────────────────────────────────────────────────────────────────

\textbf{Use case:} A transport planner wants to quantify how much
worse delays are during rush hour.

\subsubsection{SQL}

\begin{lstlisting}[style=sql, caption={Q8 --- peak vs.\ off-peak
  (\texttt{sql/queries/q8\_peak\_vs\_offpeak.sql}).}]
SELECT
    f.line_code,
    CASE
        WHEN extract(hour FROM t.poll_at_local) BETWEEN 7 AND 9
          OR extract(hour FROM t.poll_at_local) BETWEEN 17 AND 19
        THEN 'peak'
        ELSE 'offpeak'
    END                                              AS period,
    round(avg(f.mean_delay_s)::numeric,    2)       AS avg_delay_s,
    round(avg(f.mean_lateness_s)::numeric, 2)       AS avg_lateness_s,
    round(stddev(f.mean_delay_s)::numeric, 2)       AS stddev_delay_s,
    sum(f.n)                                        AS total_trains,
    count(*)                                        AS n_observations
FROM core.fact_delay_events f
JOIN core.dim_time          t ON t.poll_at_utc = f.poll_at_utc
GROUP BY f.line_code, period
ORDER BY f.line_code, period;
\end{lstlisting}

\subsubsection{Relational algebra}

The query uses an extended projection $\pi^+$ to add a computed attribute
\texttt{period} (derived from the \texttt{CASE} expression), then aggregates.
The join to \texttt{dim\_time} is needed for \texttt{poll\_at\_local}.
\begin{align*}
  &\gamma_{\texttt{line\_code},\,\texttt{period};\;
    \text{avg}(\texttt{mean\_delay\_s}),\,\text{avg}(\texttt{mean\_lateness\_s}),\,\text{sum}(\texttt{n})} \\
  &\quad\left(
    \pi^+_{\texttt{period}\,\leftarrow\,\text{CASE}\,\ldots}
    \!\left(
      \texttt{fact\_delay\_events} \bowtie_{\texttt{poll\_at\_utc}} \texttt{dim\_time}
    \right)
  \right)
\end{align*}

where $\pi^+_{c \leftarrow \text{expr}}$ denotes an extended projection adding
the computed attribute $c$; \texttt{period} $= $ ``peak'' if the local hour
falls in 07:00--09:59 or 17:00--19:59, ``offpeak'' otherwise.

\subsubsection{EXPLAIN ANALYZE}

Executed in 2,194\,ms. The planner follows these steps:
\begin{enumerate}
  \item \textbf{Parallel Seq Scan} on the fact table. \textbf{No index is used.}
    There is no \texttt{WHERE} clause to filter rows; every single row must be
    processed.
  \item \textbf{Hash Join} to \texttt{dim\_time} (37,146 rows; the hash table
    is ${\sim}$2.3\,MB, well within \texttt{work\_mem}).
  \item \textbf{Partial HashAggregate} performed by each parallel worker into
    a 10-entry hash table (5 lines $\times$ 2 periods).
  \item \textbf{Gather + Finalize HashAggregate} to merge worker outputs.
\end{enumerate}
No index is used because the query has 100\% selectivity: ignoring the index
in favour of a parallel sequential scan is the mathematically optimal choice.
This is an important counterpoint to Q5/Q6/Q7 --- indexes are not always
the answer.

\subsubsection{Sample output}

\begin{table}[H]
\centering
\begin{tabular}{llrrrr}
\toprule
\textbf{Line} & \textbf{Period} & \textbf{avg\_delay\_s} & \textbf{avg\_lateness\_s}
  & \textbf{total\_trains} & \textbf{n\_obs} \\
\midrule
RER A & offpeak &  78.02 &  88.02 & 3{,}569{,}297 & 2{,}273{,}622 \\
RER A & peak    &  88.36 & 103.46 & 2{,}496{,}086 & 1{,}354{,}143 \\
RER B & offpeak & 106.55 & 113.69 & 2{,}710{,}145 & 1{,}482{,}725 \\
RER B & peak    & 125.07 & 138.74 & 1{,}984{,}514 &   933{,}971 \\
RER C & offpeak &  89.55 &  95.69 & 2{,}381{,}729 & 1{,}565{,}332 \\
RER C & peak    & 118.53 & 130.10 & 1{,}829{,}098 & 1{,}058{,}777 \\
RER D & offpeak & 108.17 & 113.26 & 1{,}816{,}177 & 1{,}235{,}480 \\
RER D & peak    & 112.09 & 120.78 & 1{,}422{,}276 &   865{,}218 \\
RER E & offpeak &  71.81 &  80.86 & 1{,}103{,}968 &   644{,}316 \\
RER E & peak    &  77.38 &  89.90 &   750{,}008 &   401{,}629 \\
\bottomrule
\end{tabular}
\caption{Q8 results (full corpus): peak delays are 10--25\,s higher across all lines.}
\label{tab:q8}
\end{table}


% ──────────────────────────────────────────────────────────────────────
\subsection{Query performance summary}
% ──────────────────────────────────────────────────────────────────────

\begin{table}[H]
\centering
\small
\begin{tabular}{llrll}
\toprule
\textbf{Query} & \textbf{Description} & \textbf{Time (ms)} &
  \textbf{Plan} & \textbf{Index} \\
\midrule
Q5 & Station timeline     & 9,539$^\ast$ & Bitmap Index Scan &
  \texttt{idx\_fact\_stop} \\
Q6 & Line comparison (7d) & 4,323 & Bitmap Index Scan &
  \texttt{idx\_fact\_line\_time} \\
Q7 & Geographic map (1mo) & 4,516 & Bitmap Index Scan &
  \texttt{idx\_fact\_line\_time} \\
Q8 & Peak vs off-peak     & 2,194 & Parallel Seq Scan & --- \\
\bottomrule
\end{tabular}
\caption{Query performance summary. $^\ast$Cold cache; warm cache $<$1\,s.}
\label{tab:query-summary}
\end{table}

These four queries collectively illustrate the PostgreSQL planner's sophisticated
cost-based logic. Q5 demonstrates the sheer speed of point lookups via indexes
for highly selective queries. Q6 highlights the planner's adaptability, using
the same index or discarding it entirely depending on the date range width.
Q7 proves the efficacy of carefully ordered compound indexes. Finally, Q8 serves
as a vital counterpoint: when the entire dataset is needed, ignoring indexes in
favour of sequential scans is not a failure but the optimal mathematical path.


\newpage
% =============================================================================
%  SECTION 8 — OPTIMISATION ANALYSIS                                 [N]
% =============================================================================
\section{Optimisation analysis}
\label{sec:optimisation}

\subsection{The forced-plan experiment}

To validate the mathematical decisions of the PostgreSQL query planner, we
conducted a controlled degradation experiment on a full-table aggregation query.
We ran a weekly aggregation touching all 11.8\,million rows in two ways:
\begin{enumerate}
  \item \textbf{Natural plan:} The planner was allowed to choose freely,
    opting for a sequential scan.
  \item \textbf{Forced index plan:} We disabled sequential scans by executing
    \texttt{SET enable\_seqscan = off}, forcing the engine to use the B-tree index.
\end{enumerate}
The results confirm the extreme accuracy of the planner's cost model.

\begin{table}[H]
\centering
\begin{tabular}{lrrr}
\toprule
\textbf{Plan} & \textbf{Execution time} & \textbf{Buffer pages} &
  \textbf{I/O pattern} \\
\midrule
Natural (Seq Scan)  & 1,543\,ms   & 146,000      & Sequential \\
Forced (Index Scan) & 445,588\,ms & 11,800,000   & Random \\
\bottomrule
\end{tabular}
\caption{Forced-plan comparison: $289\times$ slowdown when the index is forced on a full-table scan.}
\label{tab:forced-plan}
\end{table}

\begin{figure}[H]
  \centering
  \includegraphics[width=0.72\textwidth]{figures/forced_plan_bars.pdf}
  \caption{Execution time comparison between the natural plan and the forced index plan
    (log scale). The $289\times$ slowdown confirms the planner's cost model.}
  \label{fig:forced-plan-bars}
\end{figure}

\subsection{Why the forced plan is slower}

The stark performance difference (a $289\times$ slowdown) is rooted entirely
in hardware I/O mechanics. A sequential scan reads the table's ${\sim}$130,000
data pages precisely once, in continuous physical order on disk. The OS
read-ahead prefetcher detects this linear pattern and preemptively loads
upcoming pages into RAM, making sequential I/O extremely fast.

Conversely, the forced index plan must traverse the B-tree to find 11.8\,million
index entries. Because the data on disk is not physically sorted by the index key,
each entry points to a row scattered randomly across the 130,000 pages. This
forces the disk to perform \emph{random I/O}. The buffer statistics reveal the
cost: the same physical pages are accessed over and over out of order, resulting
in 11.8\,million buffer accesses instead of 146,000. Considering that random I/O
is roughly $100\times$ more expensive per page on HDDs and ${\sim}10\times$ on
SSDs compared to sequential reading, it is clear that for full-table aggregations
the index adds catastrophic overhead without reducing the amount of data processed.

\subsection{The selectivity threshold and cost model}

The results from Q5, Q6, Q7, Q8, and the forced-plan experiment unify under a
single principle: cost-based optimisation. The planner evaluates plans in
abstract ``cost units'' proportional to disk I/O:
\begin{itemize}
  \item Index scan cost: $\approx (\text{index traversal}) + (\text{random page reads})$
  \item Seq scan cost: $\approx (\text{sequential page reads}) / (\text{parallelism factor})$
\end{itemize}
The break-even point (selectivity threshold) occurs where these two costs
intersect, typically around 10--15\%. Table~\ref{tab:selectivity_threshold}
maps the crossover behaviour across our query catalogue.

\begin{table}[H]
\centering
\begin{tabular}{lrl}
\toprule
\textbf{Selectivity} & \textbf{Query} & \textbf{Planner decision} \\
\midrule
0.7\%  & Q5 (station point lookup)   & Index Scan wins definitively. \\
4.7\%  & Q6 (7-day window)           & Index Scan wins (below threshold). \\
100\%  & Q8 (full corpus)            & Seq Scan wins; index correctly ignored. \\
100\%  & Forced-plan experiment      & Forced Index Scan: $289\times$ slower. \\
\bottomrule
\end{tabular}
\caption{The relationship between query selectivity and planner decisions.}
\label{tab:selectivity_threshold}
\end{table}

\begin{figure}[H]
  \centering
  \includegraphics[width=0.88\textwidth]{figures/selectivity_crossover.pdf}
  \caption{Conceptual selectivity--cost crossover: below ${\sim}$10--15\%
    the index plan is cheaper; above that the sequential scan wins.
    Points mark the four queries in our catalogue.}
  \label{fig:selectivity-crossover}
\end{figure}

\subsection{Implications for schema and application design}

This analysis dictates several practical design principles for the RER database
and dashboard application:
\begin{itemize}
  \item \textbf{Workload-driven indexing:} Index design must be driven by actual
    query patterns, not instinct. Each secondary index occupies ${\sim}$81\,MB
    on disk and slows down bulk loads. This is why we created all indexes
    \emph{after} the raw data load was complete.
  \item \textbf{Application-level caching:} For analytical queries with low
    selectivity (such as Q8), hitting the database repeatedly is expensive.
    The Streamlit dashboard mitigates this via \texttt{@st.cache\_data}, ensuring
    that repeated queries with identical parameters skip the database round-trip
    entirely.
  \item \textbf{Memory tuning:} Setting \texttt{work\_mem = '64MB'} in the
    database session eliminates disk spills during complex join-heavy operations
    (Q5, Q7), allowing intermediate hash tables and sort buffers to remain
    entirely in RAM.
\end{itemize}


\newpage
% =============================================================================
%  SECTION 9 — DASHBOARD                                             [L]
% =============================================================================
\section{Dashboard}
\label{sec:dashboard}

The final deliverable of this project is an interactive web-based dashboard
that allows users to explore the RER delay dataset visually while observing
the real-time performance of the underlying PostgreSQL engine.

\subsection{Architecture}

The application is built using the \textbf{Streamlit} framework, which allows
rapid development of data-driven web interfaces in Python. The technology stack
is:
\begin{itemize}
  \item \textbf{psycopg}: PostgreSQL driver for managing database connections.
  \item \textbf{pandas}: Data manipulation and intermediate format between SQL
    results and visualisations.
  \item \textbf{Plotly Express}: Interactive line and bar charts.
  \item \textbf{Folium} and \textbf{streamlit-folium}: Interactive geographic
    station map rendered via Leaflet.js.
\end{itemize}

The project structure supports Streamlit's multi-page functionality:
\begin{itemize}
  \item \texttt{dashboard/app.py}: Landing page and entry point (Line Comparison, Q6).
  \item \texttt{dashboard/db.py}: Shared module containing the database connection
    logic and the official IDFM line colour palette.
  \item \texttt{dashboard/pages/}: One Python file per additional page
    (Station Explorer, Peak Analysis, Station Timeline).
\end{itemize}

Streamlit reruns the entire script on every user interaction. Performance is
maintained through two caching layers: \texttt{@st.cache\_resource} persists
the database connection across reruns (avoiding connection-open overhead), and
\texttt{@st.cache\_data} caches query results keyed on their parameters
(avoiding redundant database round-trips for repeated identical queries).

\begin{figure}[H]
  \centering
  \input{figures/dashboard_architecture.tex}
  \caption{Dashboard architecture: Streamlit application with two caching layers
    sitting above the PostgreSQL database.}
  \label{fig:dashboard-arch}
\end{figure}

\subsection{Database connection}

The \texttt{db.py} module implements a centralised connection manager.
The \texttt{@st.cache\_resource} decorator ensures the application maintains
a single connection pool, preventing the overhead of re-opening connections on
every widget change. Query results are additionally cached with
\texttt{@st.cache\_data} keyed on user-selected parameters, so repeated
interactions do not re-hit the database.

To ensure security, the dashboard strictly uses \emph{parameterised queries}
(psycopg's \texttt{\%(name)s} syntax) rather than Python f-strings, effectively
eliminating SQL injection risks. The connection is also configured with
\texttt{SET work\_mem = '64MB'} to optimise the complex Hash Joins required
for station-level aggregations (Q5, Q7).

\subsection{Page 1: Line comparison (Q6)}

The Line Comparison page provides a high-level overview of the five RER lines.
It features a date range picker allowing users to filter the 5-month corpus,
and an interactive horizontal bar chart showing average lateness per line using
official IDFM line colours. This page serves as a primary demonstration of the
query planner's adaptive behaviour: a narrow date range (e.g.\ 7 days) triggers
a \textbf{Bitmap Index Scan} on \texttt{idx\_fact\_line\_time}, whereas selecting
the full corpus causes the planner to switch to a \textbf{Parallel Sequential Scan}
for better throughput.

\begin{figure}[H]
  \centering
  \includegraphics[width=0.9\textwidth]{screenshots/page1_line_comparison.jpeg}
  \caption{Dashboard Page 1: line comparison bar chart (Q6, 7-day window).}
  \label{fig:page1}
\end{figure}

\subsection{Page 2: Station explorer (Q7 + Q5)}

The Station Explorer provides a geographic view of the network's performance.
Users select a line and date range to view a Folium map populated with circle
markers colour-coded from green (low delay) to red (high delay) and sized by
lateness severity. An integrated drill-down feature lets users select a specific
station from the map to trigger a Q5 timeline directly below. This page
illustrates the compound index efficiency of \texttt{idx\_fact\_line\_time}
for the geographic filter (Q7) and the point-lookup speed of
\texttt{idx\_fact\_stop} for the timeline (Q5, 0.7\% selectivity).

\begin{figure}[H]
  \centering
  \includegraphics[width=0.75\textwidth]{screenshots/page2_station_explorer1.jpeg}
\end{figure}

\begin{figure}[H]
  \centering
  \includegraphics[width=0.9\textwidth]{screenshots/page2_station_explorer2.jpeg}
  \caption{Dashboard Page 2: geographic map (Q7) with Q5 drill-down timeline.}
  \label{fig:page2}
\end{figure}

\subsection{Page 3: Peak vs.\ off-peak (Q8)}

This page visualises the impact of rush hours on train punctuality. It processes
the full row corpus to produce a grouped bar chart comparing peak periods
(07:00--09:59 and 17:00--19:59) against off-peak performance across all five
lines. The dashboard highlights a key finding: peak delays are consistently
10--30\,s higher across all lines. Database analysis confirms that a
\textbf{Parallel Sequential Scan} is the optimal strategy here, as 100\% of
the rows must be categorised via the SQL \texttt{CASE} expression.

\begin{figure}[H]
  \centering
  \includegraphics[width=0.9\textwidth]{screenshots/page3_peak_offpeak.jpeg}
  \caption{Dashboard Page 3: peak vs.\ off-peak grouped bar chart (Q8).}
  \label{fig:page3}
\end{figure}

\subsection{Page 4: Station timeline (Q5 standalone)}

The standalone timeline page is designed for focused auditing of any of the
475 stations in the network. It includes a searchable dropdown for station
selection and a line chart displaying both \texttt{avg\_delay\_s} and
\texttt{avg\_lateness\_s} over time, along with summary metric cards and a
raw data table. Due to the B-tree structure of \texttt{idx\_fact\_stop},
these point lookups --- filtering for approximately 0.7\% of the data ---
are near-instantaneous, typically executing in under 200\,ms on a warm cache.

\begin{figure}[H]
  \centering
  \includegraphics[width=0.9\textwidth]{screenshots/page4_station_timeline.jpeg}
  \caption{Dashboard Page 4: station delay timeline (Q5).}
  \label{fig:page4}
\end{figure}

\subsection{Query--index--plan correspondence}

Table~\ref{tab:dashboard-index-map} summarises the mapping between each
dashboard page, its underlying query, the index used, and the execution plan
chosen by the planner. This table ties the dashboard section directly to the
indexing and optimisation analysis of Sections~\ref{sec:indexing}
and~\ref{sec:optimisation}.

\begin{table}[H]
\centering
\small
\begin{tabular}{lllll}
\toprule
\textbf{Page} & \textbf{Query} & \textbf{Index used} &
  \textbf{Plan} & \textbf{Selectivity} \\
\midrule
1 Line comparison  & Q6 & \texttt{idx\_fact\_line\_time} &
  Bitmap / Seq (adaptive) & 4.7\% / 100\% \\
2 Station explorer & Q7 & \texttt{idx\_fact\_line\_time} &
  Bitmap Index Scan & ${\sim}$8\% \\
2 Drill-down       & Q5 & \texttt{idx\_fact\_stop}       &
  Bitmap Index Scan & 0.7\% \\
3 Peak vs off-peak & Q8 & ---                            &
  Parallel Seq Scan & 100\% \\
4 Station timeline & Q5 & \texttt{idx\_fact\_stop}       &
  Bitmap Index Scan & 0.7\% \\
\bottomrule
\end{tabular}
\caption{Dashboard pages and their corresponding queries, indexes, and execution plans.}
\label{tab:dashboard-index-map}
\end{table}

\subsection{User stories}

\begin{table}[H]
\centering
\small
\begin{tabular}{p{3.8cm}p{4cm}p{5cm}}
\toprule
\textbf{User} & \textbf{Action} & \textbf{Database response} \\
\midrule
Commuter checking last week &
  Opens page 1, selects a 7-day window &
  Planner uses \texttt{idx\_fact\_line\_time} (4.7\% selectivity). \\
\addlinespace
Operator auditing a station &
  Opens page 4, selects station &
  Planner uses \texttt{idx\_fact\_stop} (0.7\% selectivity, $<$1\,s warm). \\
\addlinespace
Journalist mapping delays &
  Opens page 2, selects RER A + month &
  Planner uses compound index on line + date range. \\
\addlinespace
Researcher checking peak hours &
  Opens page 3 &
  Full-table seq scan (100\% selectivity, index not used). \\
\bottomrule
\end{tabular}
\caption{User stories and database behaviour for each dashboard page.}
\label{tab:user-stories}
\end{table}

\newpage

% =============================================================================
%  SECTION 10 — CONCLUSIONS                                          [All]
% =============================================================================
\section{Conclusions}
\label{sec:conclusions}

\subsection{Summary of findings}

This project demonstrated the lifecycle of a relational database system applied to a real-world application.
The main findings are:

\begin{itemize}
  \item \textbf{Data quality.}
    Automated profiling confirmed that the our personal IDFM delay dataset is clean and
    internally consistent.

  \item \textbf{Schema design.}
    The star-snowflake schema correctly models the IDFM stop hierarchy while
    avoiding double-counting in station-level aggregations.
    The three-level chain (\texttt{dim\_stop} $\to$ \texttt{dim\_monomodal\_stop}
    $\to$ \texttt{dim\_station}) is the key design insight of the project, arising
    directly from data profiling.

  \item \textbf{ETL reliability.}
    The idempotent pipeline reliably loads 11.8\,M rows and 2,533 station
    reference rows.
    Bulk \texttt{COPY} proved 10--100$\times$ faster than row-by-row \texttt{INSERT},
    and deferring index creation to after the bulk load further reduced load time.

  \item \textbf{Index effectiveness.}
    B-tree indexes reduce query times dramatically for selective queries:
    Q5 (0.7\% selectivity, point lookup) executes in under 200\,ms on a warm cache,
    compared to a potential 9.5\,s cold-cache scan without the index.

  \item \textbf{Query planner accuracy.}
    The forced-plan experiment produced a $289\times$ slowdown by forcing an index
    scan on a 100\%-selectivity query.
    This confirms that the planner's cost-based decision to use a parallel
    sequential scan for Q8 is not just correct but mathematically optimal: random
    I/O for 11.8\,M scattered pages is catastrophically expensive.

\end{itemize}

\subsection{Limitations}

\begin{itemize}

  \item \textbf{Single-machine performance.}
    All query execution times were measured on a development laptop.
    The absolute values (milliseconds) are hardware-dependent and should not be
    compared to production server benchmarks.
    The \emph{relative} comparisons (index vs.\ seq scan, peak vs.\ off-peak) are
    hardware-independent.

  \item \textbf{Temporal coverage.}
    The corpus covers approximately five months (November 2025 -- April 2026).
    This window includes winter and spring but excludes summer holiday traffic
    patterns, which are known to differ significantly on the RER network.
\end{itemize}

\subsection{Future work}

\begin{itemize}
  \item \textbf{Table partitioning.}
    Partitioning \texttt{fact\_delay\_events} by month (using PostgreSQL declarative
    partitioning on \texttt{poll\_at\_utc}) would allow date-range queries to prune
    entire partitions without relying on B-tree indexes, potentially reducing
    query time for medium-selectivity queries below the current crossover threshold.

  \item \textbf{Materialised views.}
    Pre-computing weekly or monthly aggregations per line and per station as
    materialised views would make dashboard interactions near-instantaneous for
    coarse-grained analyses, at the cost of incremental refresh overhead on new data.

  \item \textbf{Extended time series.}
    Adding more daily files (covering summer 2026 and beyond) would allow
    seasonal pattern analysis and validate whether the RER\,B/C lateness
    observed in this corpus is structural or seasonal.

  \item \textbf{Remote deployment.}
    Deploying the Streamlit dashboard on a shared server (e.g.\ Streamlit Community
    Cloud) with a hosted PostgreSQL instance would allow the entire team and
    anyone in general to access the application without a local installation.

\end{itemize}

% =============================================================================
%  APPENDIX
% =============================================================================
\appendix

\section{Full DDL listing}
\label{app:ddl}

The complete DDL for the \texttt{core} schema is shown inline in
Section~\ref{sec:schema}. The staging DDL is in
\texttt{sql/schema/01\_staging.sql}; the full core DDL is in
\texttt{sql/schema/02\_core.sql}.

\section{Full SQL for all 4 queries}
\label{app:sql}

The complete, annotated SQL for Q5--Q8 is shown inline in
Section~\ref{sec:queries}.
The original files are in \texttt{sql/queries/} and include, among other queries, 
parameter descriptions and usage examples as comments.

\end{document}
